% Sumber: docs/proposal-latex/references.bib (entri terverifikasi), docs/proposal/research/STATE_OF_ART_AND_GAP.md.
\chapter{Tinjauan Pustaka}
\label{chap:literature-review}

\section{Pengenalan Ekspresi Wajah dan Dataset Acuannya}

Sistem FER yang dibangun di atas arsitektur konvolusional atau arsitektur ringan untuk perangkat mobile hampir selalu dilatih dan diuji pada salah satu dari segelintir dataset acuan (FER2013, AffectNet, RAF-DB), dan survei \citet{li2020deep} atas bidang ini memperlakukan kondisi tersebut sebagai hal yang lumrah, bukan pengecualian. HSEmotion \citep{savchenko2022hsemotion}, model yang dievaluasi di sepanjang tesis ini, mengikuti pola tersebut: didistribusikan sebagai klasifier terlatih berbasis EfficientNet yang dioptimalkan untuk kecepatan inferensi, dilatih pada data keluarga AffectNet \citep{mollahosseini2017affectnet}, dan dipublikasikan sebagai artefak perangkat lunak (rilis pustaka disertai makalah teknis singkat), bukan sebagai kontribusi metodologis bagi penelitian FER itu sendiri. Perannya dalam tesis ini sengaja dibatasi: dievaluasi sebagai klasifier produksi yang sudah jadi dan tetap, tidak dilatih ulang atau disesuaikan (\textit{fine-tuned}), dan yang menjadi perhatian tesis ini adalah bagaimana klasifier tetap tersebut berperilaku di luar data tempat ia dibangun, bukan untuk memajukan teknik pemodelan FER.

\citet{li2020deeper} adalah pustaka terdahulu yang paling relevan secara langsung di sini. Analisis mereka atas bias dataset pada benchmark FER menunjukkan bahwa angka akurasi rata-rata suatu model bisa menutupi seberapa timpang akurasi itu sebenarnya tersebar di antara subkelompok dalam populasi dataset itu sendiri, dan ini adalah sifat dari cara benchmark tersebut disusun, bukan kebetulan. RQ1 tesis ini (Bab~\ref{chap:introduction} \S\ref{sec:rq}) pada dasarnya adalah versi satu-model-satu-sampel dari pertanyaan yang sama, diajukan pada data yang sama sekali tidak dicakup oleh literatur benchmark tersebut: bukan ``apakah AffectNet timpang secara internal,'' melainkan ``apa yang terjadi pada model yang sebagian dilatih dari data keluarga AffectNet ketika distribusi masukannya bergeser sepenuhnya ke luar itu.''

\section{Representasi Wajah dan Peran ArcFace}

ArcFace \citep{deng2019arcface} adalah model \emph{verifikasi} wajah: fungsi \textit{loss} \textit{additive angular margin}-nya dirancang untuk merapatkan \textit{embedding} dari identitas yang sama dan merenggangkan \textit{embedding} dari identitas yang berbeda, dan keluaran 512-dimensinya dibangun untuk pencocokan identitas, bukan untuk membedakan ekspresi. Memakai \textit{embedding} ArcFace yang dibekukan sebagai masukan bagi klasifier yang dilatih terpisah, konfigurasi yang dievaluasi pada EXP-007 (Bab~\ref{chap:methodology} \S\ref{sec:arcface-config}), adalah strategi ekstraksi fitur: representasi yang sudah dilatih untuk identitas itu dipakai persis seperti aslinya, dan hanya sebuah kepala klasifikasi Regresi Logistik yang ringan yang dilatih di atasnya memakai label milik proyek ini sendiri. Hal ini perlu dinyatakan dengan tegas, karena pemahaman yang keliru, bahwa ArcFace entah bagaimana dilatih ulang atau disesuaikan untuk pengenalan ekspresi dalam proyek ini, tidak didukung oleh bukti yang ada (Bab~\ref{chap:methodology} \S\ref{sec:arcface-config} mendokumentasikan secara persis bobot mana yang dibekukan dan bobot mana yang dilatih).

Strategi ini sendiri bukan hal baru. \citet{altaha2023facial} memadukan fitur ArcFace dengan jaringan TinySiamese untuk pengenalan ekspresi, dan \citet{waldner2024pairwise} memakai fitur ArcFace langsung untuk membedakan pasangan ekspresi pada AffectNet. Keduanya terbit sebelum perbandingan EXP-007 dalam proyek ini. \textbf{Tesis ini tidak mengklaim bahwa memadukan fitur ArcFace dengan klasifier lanjutan adalah teknik baru}; apa pun kontribusi tesis ini melalui RQ3 terletak pada perbandingan retrospektif yang spesifik terhadap garis dasar terlatih pada sampel di luar distribusi data latih, bukan pada gagasan ekstraksi fitur itu sendiri. Memperlakukan implementasi dari teknik yang sudah dipublikasikan sebagai kontribusi ilmiah tersendiri akan melebih-lebihkan apa yang sebenarnya dibuktikan oleh EXP-007, dan tesis ini menghindari pembingkaian semacam itu.

\section{Praproses Kontras dan CLAHE}

Penyetaraan histogram adaptif semacam yang diimplementasikan di sini, dibatasi klipnya (\textit{clip-limited}), diterapkan pada kanal kecerahan, dengan batas klip yang berubah-ubah mengikuti estimasi kecerahan tiap potongan, adalah teknik normalisasi pencahayaan yang umum dalam \textit{pipeline} praproses citra; tidak ada satu sumber baku yang menentukan parameterisasi yang ``benar'' secara universal, dan rentang batas klip serta skala berbasis-kecerahan yang dipakai pada EXP-006 (Bab~\ref{chap:methodology} \S\ref{sec:clahe-config}) adalah pilihan konfigurasi proyek ini sendiri, bukan nilai yang diambil dari rujukan tertentu. \TODO{Rujukan metodologi CLAHE yang primer (misalnya rumusan penyetaraan histogram adaptif asli oleh Zuiderveld) akan memperkuat bagian ini apabila dapat diverifikasi secara independen; belum ada yang terkonfirmasi dalam daftar pustaka proyek ini per revisi ini.} Yang relevan untuk RQ2 lebih sempit daripada sekadar apakah CLAHE ``berhasil'' secara umum: EXP-006 menguji satu konfigurasi tertentu, satu kali saja, tanpa pengulangan independen atau uji signifikansi atas perubahan akurasi yang teramati, dan Bab~\ref{chap:results} \S\ref{sec:rq2-results} menegaskan secara eksplisit bahwa hasil tersebut dibaca sebatas lingkup itu, tidak lebih luas.

\section{Evaluasi Bertingkat-Subkelompok dan Batasannya}

Kontribusi \citet{buolamwini2018gender} pada literatur evaluasi keadilan bersifat metodologis, bukan semata-mata empiris: menstratifikasi evaluasi berdasarkan variabel subkelompok memunculkan ketimpangan yang akan hilang ditelan rata-rata agregat, dan melakukan hal itu dengan sendirinya tidak mensyaratkan, atau mengklaim membuktikan, penjelasan sebab-akibat mengapa ketimpangan itu ada. Perbedaan inilah yang paling banyak dipakai tesis ini pada Bab~\ref{chap:results} ketika melaporkan selisih akurasi antar-kelompok kecerahan kulit untuk RQ1: selisih tersebut dilaporkan sebagai statistik deskriptif yang distratifikasi, disertai uji asosiasi statistik yang eksplisit (Bab~\ref{chap:methodology} \S\ref{sec:metrics}), bukan sebagai bukti mekanisme sebab-akibat. \citet{chhua2024bias} menerapkan kerangka ketimpangan subkelompok yang sebanding, khusus untuk FER, meskipun pada model fondasi multimodal besar, bukan pada klasifier ringan dan khusus seperti HSEmotion. Arsitekturnya cukup berbeda sehingga temuan mereka tidak bisa langsung diterapkan di sini, tetapi kerangka berpikir mereka (memperlakukan selisih akurasi antar-subkelompok sebagai objek yang diteliti, terpisah dari klaim sebab-akibat apa pun) adalah preseden metodologis terdekat yang ditemukan untuk analisis proksi warna kulit dalam tesis ini.

Perlu juga dinyatakan dengan jelas apa yang \emph{bukan} kelompok warna kulit yang dipakai di seluruh tesis ini: bukan klasifikasi tipe kulit skala Fitzpatrick, bukan ukuran demografis yang dilaporkan sendiri atau tervalidasi secara eksternal, dan bukan label etnis. Kelompok ini adalah statistik kecerahan pada ruang warna LAB yang dihitung langsung dari piksel potongan wajah (Bab~\ref{chap:methodology} \S\ref{sec:dataset-overview}), dipakai di sini sebagai proksi karena tidak ada anotasi demografis untuk dataset ini dan tidak ada yang dikumpulkan. Hasil yang distratifikasi berdasarkan proksi ini menjelaskan perbedaan antara dua kelompok kecerahan dalam satu sampel; hasil ini bukan audit keadilan terhadap kategori demografis resmi mana pun, dan tesis ini tidak memakai istilah ``bias,'' ``keadilan,'' atau ``ketimpangan demografis'' sebagai kesimpulan yang sudah terbukti di mana pun dalam Bab~\ref{chap:results}.

\section{Titik-Fitur Wajah: Dibatasi untuk Jalur Pendukung Tesis Ini}

Deteksi titik-fitur wajah (\textit{facial landmark}) dibahas di sini hanya sejauh dibutuhkan untuk menjelaskan jalur pendukung Papuan V2 (Bab~\ref{chap:methodology} \S\ref{sec:methodology-landmarks}), sejalan dengan fakta bahwa titik-fitur bukan subjek penelitian tesis ini (\S\ref{sec:background}). Ada dua pemakaian yang berbeda dalam proyek ini: \textit{mesh} 468-titik dari MediaPipe, dipakai dalam \textit{pipeline} lama untuk menurunkan fitur geometris (jarak turunnya alis, bukaan mulut) yang tidak menjadi masukan bagi satu pun dari tiga eksperimen FER yang dilaporkan pada Bab~\ref{chap:results}; dan uji coba anotasi 68-titik yang terpisah dan lebih baru (konvensi ibug-300W, usulan otomatis dari submodel \texttt{landmark\_3d\_68} milik InsightFace) yang, pada saat tesis ini ditulis, belum diperiksa oleh anotator manusia. Belum ada tinjauan literatur tentang metodologi anotasi titik-fitur manual atau kesepakatan antar-anotator untuk penempatan titik-fitur dalam tesis ini, karena uji coba tersebut belum mencapai tahap di mana tinjauan semacam itu akan memberi informasi yang konkret. Tinjauan ini layak ditambahkan apabila dan ketika jalur tersebut melangkah lebih jauh dari sekadar pembuatan usulan otomatis.

\section{Posisi Tesis Ini}

Secara keseluruhan, literatur yang ada menunjukkan bahwa HSEmotion dan ArcFace masing-masing sudah banyak diteliti, bahwa ArcFace sebagai pengekstrak fitur untuk FER sudah ada presedennya, dan bahwa evaluasi bertingkat-subkelompok adalah metodologi yang mapan untuk memunculkan, bukan menjelaskan, ketimpangan akurasi. Yang ditambahkan tesis ini bukan teknik baru, melainkan pengukuran yang spesifik dan dapat ditelusuri: bagaimana dua pendekatan yang sudah dipublikasikan tersebut, di bawah satu varian praproses, berperilaku pada populasi yang tidak dievaluasi oleh satu pun studi benchmark atau studi pendahulu yang dikutip. Ini adalah klaim yang jauh lebih sempit daripada ``tesis ini mengatasi bias FER,'' dan inilah klaim yang sebenarnya dibuat tesis ini.

\section*{Sitasi yang Masih Memerlukan Verifikasi}
\addcontentsline{toc}{section}{Sitasi yang Masih Memerlukan Verifikasi}

Lima sumber yang teridentifikasi selama pencarian literatur proyek ini sengaja tidak dimasukkan ke \texttt{references.bib} karena metadata bibliografinya belum diverifikasi secara independen: Hosseini, Pourramezan Fard \& Mahoor, ``Faces of Fairness'' (arXiv:2502.11049); ``Racial Bias within Face Recognition: A Survey'' (arXiv:2305.00817); sumber primer skala tipe-kulit Fitzpatrick tahun 1988; makalah dataset primer RAF-DB (umumnya dirujuk sebagai Li, Deng \& Du, CVPR 2017); dan satu makalah MDPI tentang \textit{transfer learning} untuk FER yang baru diketahui judulnya saja. Tidak satu pun dari kelimanya dikutip di atas. Kelimanya perlu diverifikasi terhadap sumber primer (halaman penerbit, DOI, atau abstrak arXiv) sebelum revisi mendatang mengutipnya. Kelimanya dicantumkan di sini sebagai petunjuk awal, bukan sebagai klaim yang dijadikan sandaran tesis ini.
